% ============================================================
\chapter{Ethics, Bias, and Safety}
\label{ch:ethics}
% ============================================================

\begin{intuition}[title=\textbf{Responsibility in NLP}]
Language models are deployed at massive scale and influence consequential
decisions: hiring, justice, healthcare, education. They reproduce and
amplify biases present in their training data. Understanding, measuring,
and mitigating these biases is a fundamental responsibility of the NLP
practitioner.
\end{intuition}

% ------------------------------------------------------------
\section{Bias in Word Embeddings}
\label{sec:ethics:embeddings}
% ------------------------------------------------------------

\begin{definition}[Embedding Bias]
\index{bias!embeddings}
Word embeddings capture the statistical regularities of the corpus,
including social stereotypes. For example, in Word2Vec trained on Google
News:
\[
    \vec{\text{man}} - \vec{\text{woman}} \approx
    \vec{\text{computer programmer}} - \vec{\text{homemaker}}
\]
\end{definition}

\begin{theorem}[WEAT --- Word Embedding Association Test]
\index{WEAT}
The WEAT test (Caliskan et al., 2017) measures the association between
two sets of target words $X, Y$ and two sets of attribute words $A, B$:
\[
    d(X, Y, A, B) = \sum_{x \in X} s(x, A, B) - \sum_{y \in Y} s(y, A, B)
\]
where $s(w, A, B) = \frac{1}{\abs{A}} \sum_{a \in A} \cos(w, a)
- \frac{1}{\abs{B}} \sum_{b \in B} \cos(w, b)$.
Statistical significance is assessed via a permutation test.
\end{theorem}

\begin{example}[Gender Bias]
Applying WEAT with $X = \{\text{math, science, engineering}\}$,
$Y = \{\text{arts, literature, humanities}\}$,
$A = \text{male names}$, $B = \text{female names}$, reveals a
significant association between science and male, reproducing a gender
stereotype.
\end{example}

% ------------------------------------------------------------
\section{Bias in Language Models}
\label{sec:ethics:lm_bias}
% ------------------------------------------------------------

\begin{definition}[LLM Bias]
\index{bias!LLM}
LLMs exhibit bias at multiple levels:
\begin{enumerate}
    \item \textbf{Representational}: certain groups are underrepresented
    or stereotyped.
    \item \textbf{Allocative}: the model performs better for some groups
    (e.g., resume summaries favoring certain names).
    \item \textbf{Generative}: producing toxic, discriminatory, or
    inaccurate content about certain groups.
\end{enumerate}
\end{definition}

\begin{proposition}[Sources of Bias]
\begin{itemize}
    \item \textbf{Training data}: the web contains biased, sexist, and
    racist content.
    \item \textbf{Annotation}: annotators bring their own cultural and
    linguistic biases.
    \item \textbf{Training objective}: maximizing likelihood reproduces
    dominant patterns.
    \item \textbf{Deployment}: feedback biases (users interact more with
    certain content) amplify existing biases.
\end{itemize}
\end{proposition}

% ------------------------------------------------------------
\section{Fairness Metrics}
\label{sec:ethics:metrics}
% ------------------------------------------------------------

\begin{definition}[Equalized Odds]
\index{equalized odds}
A classifier satisfies \textbf{equalized odds} if its performance is
identical across all demographic groups $g$:
\[
    P(\hat{Y} = 1 | Y = y, G = g) = P(\hat{Y} = 1 | Y = y)
    \quad \forall g, \forall y \in \{0, 1\}
\]
\end{definition}

\begin{definition}[Demographic Parity]
\index{demographic parity}
\textbf{Demographic parity} requires that the decision be independent
of the group:
\[
    P(\hat{Y} = 1 | G = g) = P(\hat{Y} = 1) \quad \forall g
\]
\end{definition}

\begin{attention}[title=\textbf{Incompatibility of Fairness Criteria}]
It is generally impossible to simultaneously satisfy demographic parity,
equalized odds, and calibration (impossibility theorems of
Chouldechova, 2017, and Kleinberg et al., 2016). The choice of criterion
depends on the application context.
\end{attention}

\begin{definition}[Toxicity Bias]
\index{toxicity}
\textbf{Toxicity bias} measures a model's propensity to generate toxic
content (insults, threats, hate speech). The RealToxicityPrompts score
(Gehman et al., 2020) measures the probability of generating toxic
content from neutral prompts.
\end{definition}

% ------------------------------------------------------------
\section{Debiasing Techniques}
\label{sec:ethics:debiasing}
% ------------------------------------------------------------

\begin{definition}[Embedding Debiasing]
\index{debiasing}
The method of Bolukbasi et al.\ (2016) identifies the bias direction
$\mathbf{b}$ (e.g., the man-woman direction) and projects embeddings to
remove this component:
\[
    \mathbf{w}_{\text{debias}} = \mathbf{w} -
    (\mathbf{w} \cdot \mathbf{b}) \mathbf{b}
\]
for gender-neutral words (not definitionally gendered).
\end{definition}

\begin{remark}
Geometric embedding debiasing has been criticized: Gonen and Gold (2019)
show that biases persist in neighborhood structure even after projection.
More recent approaches intervene at the data or training level.
\end{remark}

\begin{definition}[LLM Debiasing Strategies]
\index{debiasing!LLM}
\begin{enumerate}
    \item \textbf{Data curation}: filtering toxic content, balancing
    representation.
    \item \textbf{Instruction tuning}: fine-tuning with explicit fairness
    instructions.
    \item \textbf{RLHF/DPO}: training the model to prefer unbiased
    responses via human feedback.
    \item \textbf{Prompting}: system instructions specifying equitable
    behavior.
    \item \textbf{Controlled decoding}: modifying generation
    probabilities to reduce toxicity.
\end{enumerate}
\end{definition}

% ------------------------------------------------------------
\section{Toxicity Detection}
\label{sec:ethics:toxicity}
% ------------------------------------------------------------

\begin{definition}[Toxicity Classifier]
\index{toxicity detection}
A \textbf{toxicity classifier} is a supervised model trained to detect
offensive, hateful, or dangerous content. Examples: Perspective API
(Google), OpenAI Moderation API.
\end{definition}

\begin{proposition}[Challenges of Toxicity Detection]
\begin{enumerate}
    \item \textbf{Subjective definition}: toxicity depends on context,
    culture, and intent.
    \item \textbf{Classifier bias}: toxicity models themselves can be
    biased (e.g., classifying African-American English as more toxic).
    \item \textbf{Evasion}: users circumvent filters through creative
    rephrasing.
    \item \textbf{Over-filtering}: aggressive filtering censors
    legitimate content.
\end{enumerate}
\end{proposition}

% ------------------------------------------------------------
\section{Environmental Impact}
\label{sec:ethics:environment}
% ------------------------------------------------------------

\begin{definition}[Carbon Cost of Training]
\index{environmental impact}
Training large language models consumes substantial energy. Strubell
et al.\ (2019) estimated the carbon footprint of training a Transformer
at $\sim$284 tonnes of CO$_2$ (comparable to 5 times the lifetime
emissions of an American car).
\end{definition}

\begin{keyformulas}[title=\textbf{Carbon Footprint Estimation}]
\[
    \text{CO}_2 \approx \text{PUE} \times \text{kWh} \times
    \text{carbon intensity (g CO}_2\text{/kWh)}
\]
where PUE (Power Usage Effectiveness) $\approx 1.1$--$1.4$ for modern
data centers.
\end{keyformulas}

\begin{proposition}[Reduction Strategies]
\begin{itemize}
    \item \textbf{Smaller models}: distillation, pruning, quantization.
    \item \textbf{Efficient training}: mixed precision, gradient
    checkpointing, architecture search.
    \item \textbf{Location}: train in regions with low-carbon electricity.
    \item \textbf{Reuse}: share pretrained models (HuggingFace Hub).
\end{itemize}
\end{proposition}

% ------------------------------------------------------------
\section{Responsible AI Practices}
\label{sec:ethics:practices}
% ------------------------------------------------------------

\begin{definition}[Model Cards]
\index{model cards}
A \textbf{model card} (Mitchell et al., 2019) documents a model:
intended use, limitations, known biases, evaluation metrics, training
data, and ethical considerations.
\end{definition}

\begin{definition}[Datasheets for Datasets]
\index{datasheets}
\textbf{Datasheets} (Gebru et al., 2021) document datasets: motivation,
composition, collection process, preprocessing, recommended uses, and
demographic distributions.
\end{definition}

\begin{remark}
Regulations are evolving rapidly:
\begin{itemize}
    \item \textbf{EU AI Act} (2024): risk-based classification of AI
    systems, transparency obligations.
    \item \textbf{US Executive Order on AI} (2023): safety and testing
    requirements for foundation models.
    \item Companies adopt \emph{Responsible AI frameworks} integrating
    risk assessment, regular audits, and ethics committees.
\end{itemize}
\end{remark}

% ------------------------------------------------------------
\section{Python Implementation}
\label{sec:ethics:code}
% ------------------------------------------------------------

\begin{codeblock}[title=\textbf{Bias Detection in Embeddings}]
\begin{minted}{python}
import numpy as np
from gensim.models import KeyedVectors

def weat_score(model, X, Y, A, B):
    """Compute WEAT effect size."""
    def s(w, A_set, B_set):
        a_sims = [model.similarity(w, a) for a in A_set
                  if a in model]
        b_sims = [model.similarity(w, b) for b in B_set
                  if b in model]
        return np.mean(a_sims) - np.mean(b_sims)

    x_scores = [s(x, A, B) for x in X if x in model]
    y_scores = [s(y, A, B) for y in Y if y in model]
    d = np.mean(x_scores) - np.mean(y_scores)
    std = np.std(x_scores + y_scores)
    return d / std if std > 0 else 0.0

# Example: gender bias in professions
X = ["programmer", "engineer", "scientist"]
Y = ["nurse", "teacher", "librarian"]
A = ["man", "male", "boy", "he"]
B = ["woman", "female", "girl", "she"]
# score = weat_score(model, X, Y, A, B)
\end{minted}
\end{codeblock}

\begin{codeblock}[title=\textbf{Geometric Debiasing}]
\begin{minted}{python}
def debias_embedding(embedding, bias_direction):
    """Project out the bias direction from an embedding."""
    projection = (np.dot(embedding, bias_direction)
                  * bias_direction)
    return embedding - projection

# Compute bias direction (e.g., he - she)
bias_dir = model["he"] - model["she"]
bias_dir = bias_dir / np.linalg.norm(bias_dir)

# Debias a neutral word
debiased = debias_embedding(model["programmer"], bias_dir)
\end{minted}
\end{codeblock}

% ------------------------------------------------------------
\section{Exercises}
\label{sec:ethics:exercises}
% ------------------------------------------------------------

\begin{exercise}[Bias Analysis]
Download pretrained Word2Vec or FastText embeddings. Compute the WEAT
score for gender bias in professions. Compare results between embeddings
trained on different corpora (Wikipedia vs Common Crawl).
\end{exercise}

\begin{exercise}[Debiasing]
Implement the Bolukbasi et al.\ debiasing method. Verify that biased
analogies are reduced. Test whether debiasing affects performance on a
semantic similarity task (e.g., STS-B).
\end{exercise}

\begin{exercise}[LLM Toxicity]
Use the OpenAI Moderation API (or a HuggingFace classifier) to evaluate
the toxicity of 100 generations from a language model using neutral
prompts. Which types of toxicity are most frequent?
\end{exercise}

\begin{exercise}[Model Card]
Write a complete model card for a sentiment classifier trained on
English movie reviews. Include limitations, potential biases, and usage
recommendations.
\end{exercise}
